% Part I, divide and conquer (hylomorphisms): one Python slide, then the
% switch to Haskell

\begin{frame*}[vertical-alignment = top]
  \frametitle{Divide and Conquer}
  \vspace*{4mm}
  \mainblock{qsort-direct.py}
  \codebox<2>{dl}{dr}{dr}
  \callout<2>{14.4,6.2}{divide}{dr}
  \codebox<3>{ra}{rb}{rb}
  \codebox<3>{rc}{rd}{rd}
  \callout<3>{14.4,3.8}{solve each part}{rb, rd}
  \codebox<4>{cl}{cr}{cr}
  \callout<4>{8.4,3.6}{conquer}{cl}
  \begin{onlyenv}<5->
    Factor out \code{divide} and \code{conquer}, and the recursion is
    written \textbf{once}:
    \mainblock{hylo.py}
  \end{onlyenv}
  \par
  \codebox<6>{ea}{eb}{eb}
  \codebox<6>{pa}{pb}{pb}
  \codebox<6>{pc}{pd}{pd}
  \codebox<6>{ta}{tb}{tb}
  \codebox<6>{tc}{td}{td}
  \callout<6>{14.2,1.3}{\code{tag}: which case \code{divide} found\\
    (empty: \code{None}; non-empty: the pivot \code{p}),\\ and what
    \code{conquer} needs from it}{tb, td}
  \speakernote{Quicksort, written directly \\
    Click: divide (pivot, split the rest); recursion; conquer (glue) \\
    Click: factor them out; one generic function \\
    Last click: tag = the case divide found (None: empty list; p:
    non-empty) and the data conquer needs (the pivot); conquer dispatches on
    it \\
    The pattern: a hylomorphism (divide: coalgebra, conquer: algebra)}
\end{frame*}

\begin{frame*}
  \frametitle{The Pattern, in Haskell}
  \hsblock{hylo-qsort.hs}
  \begin{itemize}
    \item<2-> Read \code{hylo} right to left: divide, solve every part
      (\code{fmap}), conquer.
  \end{itemize}
  \speakernote{From here on: Haskell; explain the syntax as needed \\
    Layer: one layer of the call tree \\
    f . g: g then f; fmap h: h on every part \\
    hylo = conquer . fmap hylo . divide: the equation Rocq proves things
    about}
\end{frame*}

\begin{frame}[vertical-alignment = top]
  \frametitle{The Call Tree of Quicksort}
  \vspace*{5mm}
  \begin{center}
    \begin{tikzpicture}[alt={The call tree of quicksort on the list 3, 1,
        4, 5, 2, drawn as a tree of layers. The root is Node 3, for the
        input 3, 1, 4, 5, 2. Its children are Node 1, for the input 1, 2,
        and Node 4, for 4, 5. Node 1 has children Leaf and Node 2; Node 4
        has children Leaf and Node 5; Node 2 and Node 5 have two Leaf
        children each. Divide builds the tree from the root down; conquer
        folds it from the leaves up.},
      level distance = 8.5mm,
      level 1/.style = {sibling distance = 52mm},
      level 2/.style = {sibling distance = 26mm},
      level 3/.style = {sibling distance = 13mm},
      every node/.style = {box, font = \small\ttfamily},
      leaf/.style = {font = \scriptsize\ttfamily, draw = black!30,
        fill = black!5},
      inp/.style = {draw = none, fill = none, font = \scriptsize,
        text = gray, left = 1.5mm}]
      \node (r) {Node 3}
        child {node (l) {Node 1}
          child {node[leaf] {Leaf}}
          child {node (l2) {Node 2}
            child {node[leaf] {Leaf}}
            child {node[leaf] {Leaf}}}}
        child {node (rr) {Node 4}
          child {node[leaf] {Leaf}}
          child {node (r2) {Node 5}
            child {node[leaf] {Leaf}}
            child {node[leaf] {Leaf}}}};
      \node[inp] at (r.west) {$[3,1,4,5,2]$};
      \node[inp] at (l.west) {$[1,2]$};
      \node[inp] at (rr.west) {$[4,5]$};
      \node[inp] at (l2.west) {$[2]$};
      \node[inp] at (r2.west) {$[5]$};
      \draw[->, very thick, accent] (-6.0, 0.2) -- node[left, draw = none,
        fill = none, note, text = accent] {\code{divide}} (-6.0, -2.8);
      \draw[->, very thick, accent2] (6.0, -2.8) -- node[right, draw = none,
        fill = none, note] {\code{conquer}} (6.0, 0.2);
    \end{tikzpicture}
  \end{center}
  \hsblock{layer.hs}
  \begin{onlyenv}<2->
    Its type is the \textbf{least fixed point} of \code{Layer}: finite
    trees, one layer per node.
    \hsblock{calltree.hs}
  \end{onlyenv}
  \par
  \begin{onlyenv}<3->
    \code{divide} builds the tree, \code{conquer} folds it; \code{hylo}
    does both at once, so the tree is \emph{virtual}: it is never built in
    memory.
  \end{onlyenv}
  \par
  \speakernote{Everyone knows this tree; the point is its type \\
    Nodes are layers: Leaf, or Node with a pivot; grey: what divide got \\
    Layer: as in qsort; x marks where the parts go \\
    Least fixed point: smallest solution of Tree = Layer Tree, finite trees
    \\
    Haskell is lazy (also infinite trees); Rocq: finite only, so divide
    must terminate (Part II) \\
    Build, then fold = hylo, with no tree (deforestation, soon) \\
    Independent subtrees: parallelism comes back later}
\end{frame}

\begin{frame}
  \frametitle{Where Hylomorphisms Come From}
  \begin{columns}
    \begin{column}{0.4\textwidth}
      \[
        \begin{tikzcd}[alt={A commuting square: h goes from A to B; divide
            goes from A to F A, one layer of sub-problems; fmap h goes from
            F A to F B, solving every sub-problem; conquer goes from F B to
            B. So h equals conquer after fmap h after divide.},
            column sep = 15mm, row sep = 13mm]
          A \arrow[r, "h", dashed] \arrow[d, "\code{divide}"'] \& B \\
          F\,A \arrow[r, "\op{fmap}\,h"'] \& F\,B
            \arrow[u, "\code{conquer}"']
        \end{tikzcd}
      \]
      \begin{center}
        \small $h = \code{conquer} \circ \op{fmap}\,h \circ
        \code{divide}$
      \end{center}
    \end{column}
    \begin{column}{0.56\textwidth}
      \small
      \renewcommand\arraystretch{1.4}
      \begin{tabular}{@{}l@{\quad}>{\raggedright\arraybackslash}p{70mm}@{}}
        1987 & \textbf{Hagino}: data types from category theory \\
        1996 & \textbf{Hu, Iwasaki, Takeichi}: found in ordinary programs
          \\
        2006 & \textbf{Capretta, Uustalu, Vene}: when $h$ is unique \\
        2020 & \textbf{Ad\'amek, Milius, Moss}: if \code{divide}
          terminates, $h$ exists and is unique \\
      \end{tabular}
    \end{column}
  \end{columns}
  \slideref{T. Hagino. \textbf{A Categorical Programming Language}. PhD thesis, 1987. \rfsep
    Z. Hu, H. Iwasaki, M. Takeichi. \textbf{Deriving Structural Hylomorphisms from Recursive Definitions}. ICFP 1996. \rfsep
    V. Capretta, T. Uustalu, V. Vene. \textbf{Recursive Coalgebras from Comonads}. Inf. Comput. 2006. \rfsep
    J. Ad\'amek, S. Milius, L. S. Moss. \textbf{On Well-Founded and Recursive Coalgebras}. FoSSaCS 2020.}
  \speakernote{F: one layer (Layer in the Haskell); divide: coalgebra,
    conquer: algebra \\
    Hagino: data types as initial algebras, final coalgebras \\
    Name: Meijer et al. 1991 (Greek matter + form: unfold, then fold) \\
    Hu et al.: the HYLO system finds hylos in ordinary code \\
    Unique h: recursive coalgebras (CUV); terminating (well-founded)
    divide suffices (AMM): key for Rocq later}
\end{frame}

\begin{frame}
  \frametitle{Why Hylomorphisms Matter}
  \begin{columns}[t]
    \begin{column}{0.31\textwidth}
      \begin{sticky}[height = 34mm]
        \textbf{General}\par
        \small Most recursive algorithms fit: folds, dynamic programming,
        mutual recursion, accumulators.
      \end{sticky}
    \end{column}
    \begin{column}{0.31\textwidth}
      \begin{sticky}[height = 34mm]
        \textbf{Laws}\par
        \small Remove passes and intermediate data. Compilers can find the
        pattern in ordinary code.
      \end{sticky}
    \end{column}
    \begin{column}{0.31\textwidth}
      \begin{sticky}[height = 34mm]
        \textbf{Parallelism}\par
        \small Independent sub-problems can run at the same time.
      \end{sticky}
    \end{column}
  \end{columns}
  \slideref{P. Wadler. \textbf{Deforestation: Transforming Programs to Eliminate Trees}. TCS 1990. \rfsep
    A. Gill, J. Launchbury, S. Peyton Jones. \textbf{A Short Cut to Deforestation}. FPCA 1993. \rfsep
    A. Takano, E. Meijer. \textbf{Shortcut Deforestation in Calculational Form}. FPCA 1995. \rfsep
    Z. Hu, H. Iwasaki, M. Takeichi. \textbf{Deriving Structural Hylomorphisms from Recursive Definitions}. ICFP 1996. \rfsep
    R. Hinze, N. Wu, J. Gibbons. \textbf{Conjugate Hylomorphisms, or: The Mother of All Structured Recursion Schemes}. POPL 2015.}
  \speakernote{General: conjugate hylomorphisms \\
    Laws: deforestation (Wadler), shortcut fusion (in GHC); HYLO system \\
    Parallel: subtrees independent (later)}
\end{frame}

\begin{frame}
  \frametitle{Fusion and Deforestation}
  \begin{center}
    \begin{tikzpicture}[alt={Fusion: a pre-processing step, then divide
        and conquer, then a post-processing step, equals divide and conquer
        with a new divide step and a new conquer step, if pre-processing
        commutes with divide and post-processing with conquer.
        Deforestation: building the call tree with ana and then
        folding it with cata equals the hylomorphism, with no tree.},
      stage/.style = {box, minimum height = 7.5mm, font = \small\ttfamily},
      lbl/.style = {font = \small\bfseries, anchor = south west,
        inner sep = 1pt},
      cond/.style = {note, font = \footnotesize, anchor = north west}]
      \node[stage, anchor = west] (a0) at (0, 0) {pre};
      \node[lbl] at ([yshift = 1mm]a0.north west) {Fusion};
      \node[stage, right = 7mm of a0] (a1) {hylo conquer divide};
      \draw[->, thick] (a0) -- (a1);
      \node[stage, right = 7mm of a1] (a2) {post};
      \draw[->, thick] (a1) -- (a2);
      \node[right = 4mm of a2, font = \Large] (a3) {$=$};
      \node[stage, right = 4mm of a3, draw = good, fill = good!8]
        {hylo conquer2 divide2};
      \node[cond] at ([yshift = -1mm]a0.south west)
        {if \code{pre} commutes with \code{divide}, and \code{post} with
        \code{conquer}};
      \node[stage, anchor = west] (c1) at (0, -2.3) {ana divide};
      \node[lbl] at ([yshift = 1mm]c1.north west) {Deforestation};
      \node[stage, right = 7mm of c1] (c2) {cata conquer};
      \draw[->, thick] (c1) -- (c2);
      \node[right = 4mm of c2, font = \Large] (c3) {$=$};
      \node[stage, right = 4mm of c3, draw = good, fill = good!8]
        {hylo conquer divide};
      \node[cond] at ([yshift = -1mm]c1.south west)
        {build the tree, then fold it $=$ no tree in memory};
    \end{tikzpicture}
  \end{center}
  \uncover<2->{All follow from \textbf{uniqueness} (\code{hylo} is the only
    solution of its equation), and are proved \textbf{once}, for every
    divide and conquer.}
  \speakernote{Fusion: post-processing moves into conquer, pre-processing
    into divide (either may be the identity) \\
    Deforestation: ana builds the tree, cata folds it; together, no tree \\
    Uniqueness again; it needs care, back in Part II}
\end{frame}

\begin{frame*}[vertical-alignment = top]
  \frametitle{Calculating Sort-Then-Double}
  \vspace*{4mm}
  \hsblock{qsort-calc.hs}
  \begin{onlyenv}<2->
    \begin{center}
      Two passes and a temporary list $\Longrightarrow$ \textbf{one pass},
      doubling each pivot.
    \end{center}
  \end{onlyenv}
  \par
  \speakernote{Spec: sort, then double \\
    Side condition: doubling distributes over concatenation \\
    Result: the program you would write by hand \\
    We will redo this in Rocq, and run it}
\end{frame*}
